\documentclass[10pt,a4paper]{article}
\usepackage[T1]{fontenc}
\usepackage{lmodern,amsmath,amssymb,amsthm,mathtools,microtype}
\usepackage[margin=24mm]{geometry}
\usepackage[colorlinks=true,linkcolor=blue,urlcolor=blue,citecolor=blue]{hyperref}
\newtheorem{theorem}{Theorem}[section]
\newtheorem{lemma}[theorem]{Lemma}
\newtheorem{proposition}[theorem]{Proposition}
\newcommand{\Q}{\mathbb Q}
\newcommand{\Z}{\mathbb Z}
\setlength{\emergencystretch}{2em}
\title{Local obstructions to small conic-chain pole systems}
\author{}\date{15 November 2026}
\begin{document}\maketitle
\begin{abstract}
A three-quadric system with four auxiliary rational coordinates
forbids every negative odd valuation of its nonzero input at every odd
prime. A related four-quadric system uses five auxiliary coordinates.
Both fail on explicit families of arbitrarily small prepared positive
inputs. More generally, every fixed rational coefficient in the
five-coordinate conic-chain family has a local obstruction on such
inputs. These results exclude replacement on the full prepared locus;
they are not lower bounds for arbitrary pole predicates or defining
quartics. A degree-four map identifies the elliptic square-class
condition for the four-coordinate system.
\end{abstract}

\section{The conic-chain predicate}
For $b\in\Q$, let $\mathcal C(b)$ be the assertion that rational
$e,p,z,P$ satisfy
\begin{equation}\label{eq:system}
 P=p^2,\qquad e^2+bP=1,\qquad z^2+bP=b.
\end{equation}
Thus $\mathcal C(b)$ uses four auxiliary coordinates and three
quadratic equations. Eliminating $P$ gives
\begin{equation}\label{eq:conics}
 e^2+bp^2=1,\qquad z^2=b(1-p^2).
\end{equation}
The sum of the squares of the three residuals in
\eqref{eq:system} is a quartic whose rational zero set expresses this
predicate. This statement alone does not give a definition of the
integers.

\begin{theorem}\label{thm:odd}
Suppose $b\in\Q^\times$ and $\mathcal C(b)$ holds. At every odd
prime $\ell$, the integer $v_\ell(b)$ is not negative and odd.
\end{theorem}
\begin{proof}
Suppose $\nu=v_\ell(b)<0$ is odd. If $p=0$, the second equation
of \eqref{eq:conics} says $z^2=b$, an immediate contradiction.
Otherwise, whenever $e^2$ and $bp^2$ are both nonzero, their valuations
have opposite parity. In the first equation no negative smaller
valuation can cancel. Consequently $v_\ell(e)\ge0$ and
$\nu+2v_\ell(p)\ge0$. The latter valuation is odd, hence positive;
it follows that $e$ is a unit and $v_\ell(p)>0$.
The same conclusion follows if $e=0$, since that case would require
$\nu+2v_\ell(p)=0$, already impossible. Now $1-p^2$ is a unit,
so $z^2=b(1-p^2)$ has odd valuation $\nu$, a contradiction.
\end{proof}

\begin{proposition}\label{prop:detector}
Let $b,w,\eta\in\Q^\times$ satisfy $wb=\eta^2$.
The equations
\[
 e^2+bp^2=1,\qquad p^2+ws^2=1
\]
are equivalent to \eqref{eq:conics} by the invertible change
$z=\eta s$. Equivalently, on putting $r=z/b$, they become
\[
 e^2+bp^2=1,\qquad p^2+br^2=1.
\]
\end{proposition}
\begin{proof}
Multiply $ws^2=1-p^2$ by $b$ and use $bw=\eta^2$.
The inverse is $s=z/\eta$. The second form follows by dividing
$z^2=b(1-p^2)$ by $b$.
\end{proof}

\section{Negative even valuations}
For an odd prime $\ell$ and $b\in\Q_\ell^\times$, write
$b=\ell^\nu c$ with $c\in\Z_\ell^\times$; the square class of
the residue of $c$ is called the unit square class of $b$.

\begin{theorem}\label{thm:even}
Let $\ell\equiv1\pmod4$ be prime. Suppose $b\in\Q_\ell^\times$
has negative even valuation and nonsquare unit square class.
Then \eqref{eq:conics} has no solution in $\Q_\ell$.
\end{theorem}
\begin{proof}
Write $v_\ell(b)=-2k$ with $k\ge1$, and write
$b=\ell^{-2k}c$ with nonsquare unit $c$. Since $-1$ is a square
in the residue field, $-c$ is also nonsquare.
If $p\ne0$ and $v_\ell(p)<k$, then $bp^2$ has negative even
valuation. The equation $e^2+bp^2=1$ would require equality of the
two negative valuations, followed by cancellation of their leading
residues. That cancellation would make $-c$ a square, impossible.
Thus $p=0$ or $v_\ell(p)\ge k$. In both cases $1-p^2$ is a unit
congruent to $1$ modulo $\ell$. The equation $z^2=b(1-p^2)$
would therefore have nonsquare unit square class $c$ on its right
side. This is impossible for a square.
\end{proof}

\begin{theorem}\label{thm:family}
For every integer $n\ge1$, put
\[
                 b_n=\frac{272}{29^{2n}}.
\]
Then $0<b_n<1$, the reduced denominator of $b_n$ is a square,
$b_n\in\Q_2^{\times2}$, and $v_2(b_n)=4$.
Nevertheless $\mathcal C(b_n)$ has no $\Q_{29}$-solution.
In particular the failure persists among prepared positive inputs
arbitrarily close to zero.
\end{theorem}
\begin{proof}
The displayed fraction is reduced, $272<29^2$, and its positive
values tend to zero. Its odd $2$-adic unit is
$17/29^{2n}\equiv1\pmod8$, so it is a $2$-adic square.
At $29$, its valuation is $-2n$ and its unit part has residue
$272\equiv11\pmod{29}$. This is a nonsquare: Euler's criterion gives
$11^{14}\equiv-1\pmod{29}$. Apply Theorem~\ref{thm:even}.
\end{proof}

Every input in Theorem~\ref{thm:family} satisfies the preparation
assumptions of the congruent-twist realization theorem~\cite{Pole}.
Consequently that theorem, with its stated arithmetic input, supplies
witnesses for the seven-coordinate pole predicate at every $b_n$,
whereas \eqref{eq:system} has no rational witness. The
four-coordinate system cannot replace that predicate on the full
prepared locus. This does not exclude a construction whose actual
inputs lie in a smaller, separately proved realization locus.

\section{The five-coordinate family}
Fix $c\in\Q$. Define $\mathcal C_c(b)$ by the four quadrics in
five auxiliary coordinates $e,p,P,D,z$:
\begin{equation}\label{eq:five}
 P=p^2,\qquad e^2+bP=1,\qquad
 D=1+(cb-1)P,\qquad z^2=bD.
\end{equation}
The coefficient $c$ is fixed, not an additional quantified coordinate.
For $c=0$, eliminating $D$ gives exactly \eqref{eq:system}.
For $c=1$, the equation defining $D$ can equally be written
$D+e^2+P=2$.

\begin{proposition}\label{prop:five-sound}
If $c\in\Q$, $b\in\Q^\times$, and $\mathcal C_c(b)$ holds,
then $v_\ell(b)$ is not negative and odd at any odd prime $\ell$
where $c$ is integral. In particular this applies to every odd prime
when $c\in\Z[1/2]$.
\end{proposition}
\begin{proof}
The first conic, exactly as in Theorem~\ref{thm:odd}, gives
$v_\ell(p)>0$ and $v_\ell(bp^2)>0$ if $v_\ell(b)$ is negative
and odd. Since $c$ is integral at $\ell$, both $cbp^2$ and $p^2$
have positive valuation. Thus $D=1+cbp^2-p^2$ is a unit, and
$z^2=bD$ has odd valuation, a contradiction. The case $p=0$ is
immediate from $z^2=b$.
\end{proof}

\begin{lemma}\label{lem:unit}
Let $\ell$ be odd, let $c\in\Z_\ell$, and suppose $c+1$ is a
unit with nonsquare residue. If $b\in\Z_\ell^\times$ has residue
$(c+1)^{-1}$, then \eqref{eq:five} has no $\Q_\ell$-solution.
\end{lemma}
\begin{proof}
The residue of $b$ is nonsquare, and $cb-1\equiv-b\pmod\ell$.
If $v_\ell(p)<0$, then $e^2=1-bp^2$ requires $-b$ to have
square unit residue. The equation
$z^2=b[1+(cb-1)p^2]$ likewise requires $-b^2$ to have square
unit residue. Their quotient would make $b$ a square residue,
a contradiction. Hence $p,e,z$ are all integral.
Reducing modulo $\ell$ gives $D=e^2$ and $z^2=be^2$.
Since $b$ is nonsquare, $e=z=0$ in the residue field. The first
conic then gives $p^2=b^{-1}$, again a contradiction.
\end{proof}

\begin{theorem}\label{thm:five-family}
For every integer $n\ge1$, put $a_n=528/31^{2n}$.
These inputs satisfy $0<a_n<1$, have reduced square denominator,
and are $2$-adic squares with $v_2(a_n)=4$.
Nevertheless $\mathcal C_1(a_n)$ has no $\Q_5$-solution.
\end{theorem}
\begin{proof}
The preparation follows from $528=16\cdot33$, $33\equiv1\pmod8$,
and $528<31^2$. For every $n$, the residue of $a_n$ at $5$ is
$3=2^{-1}$. Since $2$ is a nonsquare modulo $5$,
Lemma~\ref{lem:unit} applies with $c=1$.
\end{proof}

\section{Every fixed rational coefficient has an obstruction}

\begin{lemma}\label{lem:minus-one}
Let $\ell$ be odd and $c\in\Z_\ell$ with $c\equiv-1\pmod\ell$.
If $b$ has negative even valuation and nonsquare unit square class,
then \eqref{eq:five} has no $\Q_\ell$-solution.
\end{lemma}
\begin{proof}
Write $b=\ell^{-2k}\beta$ with $k\ge1$ and nonsquare unit $\beta$.
If $p\ne0$ and $v_\ell(p)<k$, the leading terms of the two conics
would require $-\beta$ and $c\equiv-1$ both to be squares in the
residue field. Their quotient would make $\beta$ a square.
Consequently $p=0$ or $v_\ell(p)\ge k$.
Set $p=\ell^k p_0$ and $Z=\ell^kz$. Then $p_0,e,Z$ are integral,
and reduction of the two conics gives
\[
 e^2+\beta p_0^2=1,\qquad
 Z^2=\beta(1-\beta p_0^2)=\beta e^2.
\]
Nonsquareness of $\beta$ forces $e=Z=0$ in the residue field,
leaving $\beta p_0^2=1$, impossible.
\end{proof}

\begin{lemma}\label{lem:norm}
Let $\ell$ be odd. Suppose $b$ has even valuation and nonsquare
unit square class, and suppose $d=cb-1$ has odd valuation.
Then \eqref{eq:five} has no $\Q_\ell$-solution.
\end{lemma}
\begin{proof}
For every $p$, the value $D=1+dp^2$ is either a unit congruent
to $1$ modulo $\ell$, or has negative odd valuation. Indeed,
$v_\ell(d)+2v_\ell(p)$ is odd whenever $p\ne0$, so it cannot be
zero. In the first case $bD$ has nonsquare unit square class;
in the second it has odd valuation. Neither can equal $z^2$.
\end{proof}

For example, the rational coefficient $c=-9/25$ is obstructed at
$b=272/441$. At $3$, this input has valuation $-2$ and nonsquare
unit residue $2$, whereas
\[
                 cb-1=-\frac{1497}{1225}
\]
has valuation $1$. The input is prepared: its reduced denominator is
$21^2$, its $2$-adic valuation is $4$, and its odd $2$-adic unit is
$1$ modulo $8$.

\begin{lemma}\label{lem:nonresidue}
Every nonsquare $r\in\Q^\times$ is a nonsquare unit at some odd
prime outside any prescribed finite set.
\end{lemma}
\begin{proof}
Let $\chi$ be the nontrivial quadratic character attached to
$\Q(\sqrt r)$. Choose a residue class coprime to its conductor
on which $\chi=-1$. The Chinese remainder theorem supplies a
positive representative $a$ that is also coprime to every prescribed
prime outside the conductor. Since $\chi(a)=-1$, at least one
prime factor $\ell$ of $a$ has $\chi(\ell)=-1$.
Include $2$ and all prime factors of the numerator and denominator of
$r$ among the excluded primes. Quadratic reciprocity identifies
$\chi(\ell)$ with the quadratic residue symbol of $r$ at $\ell$,
which gives the assertion. This argument uses factorization of $a$;
it does not require a theorem about primes in arithmetic progressions.
\end{proof}

\begin{lemma}\label{lem:prepared}
Fix an odd prime $\ell$.
Every prescribed unit residue modulo $\ell$ or $\ell^2$ occurs
among prepared positive rational inputs arbitrarily close to zero.
Such inputs can also have any prescribed nonsquare unit residue
and arbitrarily large negative even valuation at $\ell$.
\end{lemma}
\begin{proof}
For a prescribed unit residue $a$ modulo $\ell^j$, where $j=1$
or $2$, use the Chinese remainder theorem to choose a positive integer
$M$ with $M\equiv1\pmod8$ and $16M\equiv a\pmod{\ell^j}$.
Put $q=1+8\ell^jM$. The fractions $16M/q^{2n}$ are reduced,
have the prescribed residue, and tend to zero. They have reduced
square denominator and $2$-adic valuation $4$, and their odd
$2$-adic units are $1$ modulo $8$.
For the valuation assertion choose $M\equiv1\pmod8$ with $16M$
a nonsquare of the prescribed residue modulo $\ell$, and use
$16M/\ell^{2k}$ for sufficiently large $k$.
\end{proof}

\begin{theorem}\label{thm:all-c}
For every fixed $c\in\Q$, there are arbitrarily small positive
rational inputs with reduced square denominator that are $2$-adic
squares of $2$-adic valuation $4$, but for which $\mathcal C_c(b)$
has no solution in at least one completion of $\Q$.
\end{theorem}
\begin{proof}
The case $c=0$ is Theorem~\ref{thm:family}.

Suppose first that $c\ne0$ is not a rational square.
By Lemma~\ref{lem:nonresidue}, choose an odd prime $\ell$ at which
$c$ is a nonsquare unit. Prescribe
\[
                   b\equiv c^{-1}(1+\ell)\pmod{\ell^2}.
\]
Then $b$ is a nonsquare unit and $v_\ell(cb-1)=1$.
Lemmas~\ref{lem:norm} and \ref{lem:prepared} give the required
prepared counterinputs.

Suppose next that $c$ is a nonzero rational square but $c+1$ is
not a rational square. Choose an odd prime at which $c+1$ is a
nonsquare unit by Lemma~\ref{lem:nonresidue}, and apply
Lemmas~\ref{lem:unit} and \ref{lem:prepared}.

Finally suppose that both $c$ and $c+1$ are nonzero rational
squares. Write
\[
             c=(u/v)^2,\qquad c+1=(a/v)^2,
             \qquad u^2+v^2=a^2,
\]
with $u,v,a$ positive integers and $\gcd(u,v)=1$.
The integer $a$ is odd: if it were even, primitivity would make
$u,v$ both odd, contradicting the equation modulo $4$.
Also $a>1$, so it has an odd prime factor $\ell$.
Primitivity gives $\ell\nmid v$. Hence $c\equiv-1\pmod\ell$,
and Lemmas~\ref{lem:minus-one} and \ref{lem:prepared} finish
the proof.
\end{proof}

Theorem~\ref{thm:all-c} concerns this explicitly stated conic-chain
family and the full prepared locus. It does not exclude different
five-coordinate systems or a realization theorem restricted to the
actual inputs of a different integer construction.

\section{The elliptic square-class locus}
For $b\in\Q\setminus\{0,1\}$, consider the nonsingular elliptic
curve
\[
             E_b:\quad Y^2=X(X-b)(X-b^2).
\]
\begin{proposition}\label{prop:cover}
Every solution of \eqref{eq:conics} gives a point of $E_b$ by
\[
                 X=b^2p^2,\qquad Y=b^2pez.
\]
Away from the three points with $Y=0$, the image consists exactly
of the points satisfying
\[
 X\in\Q^{\times2},\qquad
 X-b\in-b\Q^{\times2},\qquad
 X-b^2\in-b\Q^{\times2}.
\]
Each such point has four sign-compatible preimages.
\end{proposition}
\begin{proof}
The two conics give
\[
 X-b=-be^2,\qquad X-b^2=-bz^2.
\]
Their product with $X=b^2p^2$ is $b^4p^2e^2z^2=Y^2$.
Conversely choose rational square roots of
$X/b^2$, $(b-X)/b$, and $(b^2-X)/b$. They give nonzero
$p,e,z$ satisfying both conics. Of the eight independent choices
of signs, exactly four have $b^2pez=Y$.
\end{proof}

Thus a realization theorem for the four-coordinate predicate requires
rational points in this particular simultaneous square-class locus.
Separate solvability of the two conics does not establish that
requirement; Theorem~\ref{thm:family} already obstructs it locally.

\section{Exact verification}
The accompanying program
\nolinkurl{computation/pole_parameters/conic_chain_v41/check.py}
verifies the two substitution identities, the elliptic-map identity,
the nonresidue computations at $29$ and $5$, and exact preparation
data for members of both explicit families. It also checks residue
certificates illustrating the three general local obstructions and
the rational-coefficient example $c=-9/25$. The proofs above, rather than the
finite examples, establish the statements for all primes and all
$n\ge1$.

\begin{thebibliography}{9}
\bibitem{Pole} \emph{Pole equations from simultaneous congruent-number
twists}, project companion,
\nolinkurl{papers/pole_parameters/congruent_twist_poles_20261114_v27.tex},
Theorem~2.1 and its stated arithmetic input.
\end{thebibliography}
\end{document}
